\emph{$u_n$ désigne le nombre d'individus au 1\up{er} mars $2025+n$. La population diminue
chaque année de $100$ individus ; au $1$\up{er} mars $2025$, elle compte $\np{2500}$ individus.}

\begin{enumerate}
	\item
	\begin{enumerate}
		\item L'année 2025 correspond à $n = 0$ : $u_0 = \np{2500}$.

		\item L'année 2026 correspond à $n = 1$ :
		$u_1 = u_0 - 100 = \np{2500} - 100 = \np{2400}$.

		Selon ce modèle, au 1\up{er} mars 2026, la population comptera \np{2400} individus.
	\end{enumerate}

	\item
	\begin{enumerate}
		\item Chaque année, la population diminue de 100 individus donc, pour tout entier naturel $n$ :
		
		$u_{n+1} = u_n - 100$.

		\item Pour tout entier naturel $n$, $u_{n+1} = u_n + (-100)$ : la suite $(u_n)$ est
		\textbf{arithmétique} de raison $r = -100$ et de premier terme $u_0 = \np{2500}$.

		\item Pour une suite arithmétique de premier terme $u_0$ et de raison $r$ :
		$u_n = u_0 + nr$.

		Pour tout entier naturel $n$, $u_n = \np{2500} - 100n$.
	\end{enumerate}

	\item La raison de la suite est $r = -100 < 0$. La suite $(u_n)$ est donc \textbf{décroissante}.

	\item
	\[u_n < \np{1250} \iff \np{2500} - 100n < \np{1250} \iff -100n < -\np{1250} \iff n > 12,5.\]
	Le plus petit entier $n$ qui convient, immédiatement supérieur à $12,5$,  est $n = 13$.

	C'est à partir du 1\up{er} mars $2025 + 13 = 2038$ que le nombre d'individus
	sera, selon ce modèle, inférieur à \np{1250}.
\end{enumerate}

\medskip
